\section{Computing the expected top-K count at each distance}
\label{sec:finite-distance-details}
The count $G_K(p)$ in Equation~\eqref{eq:clipped-binomial} can be evaluated directly as
\[
G_K(p)=\sum_{x=0}^{N}\min(K,x)\binom Nx p^x(1-p)^{N-x}.
\]
For large $N$, computing the binomial coefficients and powers separately can overflow or lose
small probabilities. A shorter expression follows from a counting identity. For an integer
$x\ge0$, $\min(K,x)$ counts how many thresholds $1,2,\ldots,K$ it reaches:
\[
\min(K,x)=\sum_{j=1}^{K}\ind[x\ge j].
\]
Taking expectations on both sides gives
\begin{equation}
G_K(p)=\sum_{j=1}^{K}\Pr(X\ge j),\qquad X\sim\operatorname{Binomial}(N,p).
\label{eq:clipped-survival}
\end{equation}
Use a numerically stable binomial survival function for these probabilities. The ordinary
small-example checker uses the direct sum; the simulation checker uses this alternative.
Their agreement supplies a check across two different calculations.

For the distance model, $\sum_m w_m=G_K(1)-G_K(0)=K$. Thus the weights in
Equation~\eqref{eq:finite-top-k}, divided by $K$, form a probability distribution over the
distance of a uniformly chosen true top-$K$ row. Given that row's distance, symmetry makes
its error positions uniform. The ID rule chooses among tied rows without favoring one set of
error positions. This justifies multiplying $w_m$ by the hypergeometric survival probability.

The argument would change if ties deliberately favored matching prefixes. It would also change
if we stopped candidate collection after a corpus-dependent quota. These are changes to the
selection policy, not numerical errors in the fixed-radius formula.

\par\noindent\begin{minipage}{\linewidth}
\begin{lstlisting}[caption={Expected top-K recall for a fixed prefix radius.}]
ExpectedTopKRecall(N, d, h, radius, K):
    previous_slots = 0
    expected_found = 0
    for distance in 0 .. d:
        p = F_d(distance)
        slots = sum(Pr(Binomial(N, p) >= j) for j in 1 .. K)
        rows_at_distance = slots - previous_slots
        expected_found += rows_at_distance * g_h_radius(distance)
        previous_slots = slots
    return expected_found / K
\end{lstlisting}
\end{minipage}\par

\section{A finite-sample view of prefix recall}
The population-tail approximation is not the only mathematical definition available. This
appendix shows what an exact finite-$N_*$ calculation would require, under independent,
identically distributed document score pairs and no score ties. It is useful for understanding
why knowing only a correlation is insufficient.

For a particular document with full score $s$ and prefix score $t$, classify another document
into four events:
\begin{align*}
p_{11}(s,t)&=\Pr(S>s,T>t),&p_{10}(s,t)&=\Pr(S>s,T\leq t),\\
p_{01}(s,t)&=\Pr(S\leq s,T>t),&p_{00}(s,t)&=\Pr(S\leq s,T\leq t).
\end{align*}
Among the other $N_*-1$ documents, let $(a,b,c,e)$ count these categories. A multinomial distribution counts how many independent documents land in each of the four categories, using those probabilities. The chosen document belongs to the full
top $K$ when $a+b\leq K-1$, and belongs to the prefix top $R$ when $a+c\leq R-1$.
Therefore, with $e=N_*-1-a-b-c$,
\begin{align}
P_{K,R}(s,t)=
\sum_{\substack{a,b,c\geq0;\ a+b\leq K-1;\ a+c\leq R-1\\a+b+c\leq N_*-1}}
\frac{(N_*-1)!}{a!b!c!e!},
p_{11}^{a}p_{10}^{b}p_{01}^{c}p_{00}^{e}.
\end{align}
If $F_{S,T}$ is the joint distribution of one document's scores, every document has the same event probability under the identical-sampling assumption. Summing over all $N_*$ documents and dividing by $K$ gives
\begin{equation}
\E[\operatorname{Recall@}K]=
\frac{N_*}{K}\int P_{K,R}(s,t)\,dF_{S,T}(s,t).
\label{eq:finite-recall}
\end{equation}
This assumes the candidate stage returns exactly the prefix top $R$ and the final stage uses
the same full score. Retrieving whole tied prefix keys changes that candidate rule; ties can
be handled by augmenting the ordering with document IDs or by explicitly integrating their
selection probabilities.

Equation~\eqref{eq:finite-recall} is much harder to evaluate than the conditional-normal
approximation. More importantly, it still needs $F_{S,T}$. Matryoshka training, a code length,
a cluster count and a RAM budget do not uniquely specify that distribution. Exact mathematics
can organize unknown inputs; it cannot infer them from the algorithm's name.

\section{Checklist for an honest numerical implementation}
\begin{enumerate}
\item Fix the score and tie rule. Keep original-float, truncated-float, binary-score and human
relevance references distinct.
\item Check the score identity on small vectors and enumerate the actual top $K$.
\item Check that the weighted-subset stream emits every subset once, in nondecreasing penalty,
and processes the final boundary tie before claiming exactness.
\item Count split words and leaf words separately. Use each cluster's physical width, including
padding; do not substitute the number of active documents for the bitmap length.
\item Record the ownership of codes, external IDs, posting references and key storage. Any
shared or reordered array must be stated before removing its bytes from a formula.
\item Count temporary state at its lifetime peak. A queue-capacity bound and an observed
allocation peak are different quantities.
\item Keep key attempts separate from occupied keys and returned postings. Expected first hit
is not the same threshold as finding $R$ candidates.
\item Price centroid float arithmetic separately from binary score lookups. Charge selection,
allocation, copying and gathers where the implementation performs them.
\item Use either a calibrated combined cost or a disaggregated operation model for each block.
Do not count the same memory and selection work in both.
\item Test statistical predictions on queries not used to estimate their parameters. Evaluate
recall after the actual candidate and tie rules, not merely a probability of finding some key.
\end{enumerate}

\section{Floating-point arithmetic and exactness}
The score and bound proofs use real arithmetic. A numerical implementation must ensure its
computed upper bounds remain conservative under rounding. The local C++ reference includes a
rounding allowance for bound pruning and a deterministic row-ID tie order. A report should not
claim exactness merely because an unguarded floating-point comparison resembles the algebra.

This does not require inventing a collection of hypothetical production patches. It requires
checking the numerical path used by the implementation and testing it against a full reference
scan on the same inputs. The small examples and archived candidate checks provide such evidence
for their stated scopes; they do not prove every future implementation correct.

\section{Optional expansion of the embedding cost}
The main formulas keep $T_{\rm emb}$ and $Q_{\rm emb}$ as measured model-specific functions.
If a more detailed encoder estimate is needed, its architecture must be another set of inputs,
not an unexplained constant.

For an illustrative standard dense Transformer encoder, let $n_t$ be query tokens, $L_e$ the
number of layers, $m$ hidden width, $f$ feed-forward width, and $H_a$ attention heads. For one
layer with ordinary two-matrix feed-forward computation, multiply-add contributions are
\begin{equation}
U_{\rm layer}=4n_t m^2+2n_t^2m+2n_tmf.
\end{equation}
The first term comes from the query, key, value and attention-output projections. The second
comes from attention-score and attention-value matrix products. The third comes from the
feed-forward expansion and contraction. Activation, normalization and softmax operations add
separate elementwise terms. A different block, gated feed-forward network or sparse attention
requires different counts.

A corresponding parameterized time model is
\[
T_{\rm emb}=T_{\rm tokenize}+L_e\big(U_{\rm layer}c_{\rm enc,mac}
 +U_{\rm elementwise}c_{\rm enc,elt}\big)+T_{\rm pool/project}.
\]
These effective coefficients depend on batching, device kernels, precision and memory behavior.
Alternatively use the resource bound in Equation~\eqref{eq:resource}, without adding it on top
of coefficients that already include the same work.

A non-fused attention implementation may materialize $H_a n_t^2$ attention values. With bytes
per activation $b_a$, an explicit scratch model can be written
\[
Q_{\rm emb}=b_a\big(\eta_m n_t m+\eta_f n_t f+\eta_a H_a n_t^2\big)+Q_{\rm kernel},
\]
where the $\eta$ coefficients count simultaneously live buffers, not time. Fused attention
can avoid materializing the square attention matrix, so its chosen buffer counts differ.
Layer workspaces can be reused during inference; multiplying a per-layer peak by $L_e$ is not
a general inference-RAM formula. Resident model weights are shared and counted separately.

A Matryoshka prefix normally reuses an already computed embedding. Halving the search prefix
$d$ does not automatically halve the encoder's layer count or hidden width. It reduces the
subsequent representation and search work unless the actual encoder implementation also changes.

\section{Reproduction and parameter ledger}
The editable source is \texttt{report.tex}, with chapter files in \texttt{sections/}. Mermaid
sources and rendered diagrams are in \texttt{figures/}. The build instructions are in the
adjacent \texttt{README.md}. The arithmetic checks are implemented in \texttt{verify\_math.py}.
They validate the worked examples, exact subset ordering, binomial conditional-recall example,
word widths, selected storage identities, and quoted measured rows. The statistical evidence
is preserved with its source hash.

The following ledger identifies the inputs a future calculator must obtain rather than infer
from one document count.
\small
\begin{longtable}{@{}P{.25\linewidth}P{.35\linewidth}P{.33\linewidth}@{}}
\toprule Input family & Parameters or counts & How to obtain them\\\midrule\endfirsthead
\toprule Input family & Parameters or counts & How to obtain them\\\midrule\endhead
Representation & $D,d,d_r,d_f$, preparation and score rule & fixed experimental definition; embedding/model metadata\\
Retrieval target & $K,R$, filter, reference score, tie rule & fixed task definition\\
Cluster layout & $C,n_j,\mathcal P(q),e_j$ & built index and per-query routing/filter traces\\
Storage widths & $w,b_w,b_{\rm id},b_o,b_c,b_f,b_s$ & actual types, packing and alignment\\
Baseline scoring & $g,G,H_s,c_{\rm lut}$ & source code, work counters and kernel measurements\\
Bitplane search & $L,B_{\rm node},s_j,\ell_j,f_j,V$ & configuration plus query traces; not budget alone\\
Bitplane RAM & frontier masks, current/child buffers, capacities & allocation/capacity measurements or an explicit bound\\
Key index & $I_h,h_j,U_j,\alpha,b_k,b_{\rm ref}$ & chosen lookup coordinates, index build and slot layout\\
Key search & $H_j,F_j,E_j,Q_j^{\max}$, stopping rule & actual key enumeration, posting and frontier traces\\
Prefix-trie variant & $Z_j,b_{\rm trie},u_j$, posting ranges & actual stored nodes, path and range visits\\
Hardware & compute rates, bandwidth, cache regime, SIMD, threads & calibrated for the actual workload and machine\\
Optional final stage & fetch policy, $R',u_f,d_f$, rerank score & execution pipeline and storage-residency decision\\
Statistical model & query/data distribution, covariance, score tails & declared synthetic assumptions or measured and checked estimates\\
Serving RAM & $M_{\rm resident},M_{\rm model},M_{\rm runtime},J,Q_{\rm peak}$ & ownership/lifetime analysis and process measurements\\\bottomrule
\end{longtable}
\normalsize
A numerical example can assign values to these parameters. The equations themselves remain
parameterized, so changing the workload changes the inputs instead of requiring a different
explanation of the algorithm.
